\subsection{Strategies for Enhancing Robustness and Generalizability in AI}
\label{sec:6.2.2}
The engineering here is standard, and the rest of the book assumes it exists. Adversarial training — training on inputs perturbed in a direction chosen to flip the model's prediction, which neural networks are notoriously vulnerable to — measurably increases resistance to that class of attack \autocite{szegedy2013intriguing}, \autocite{goodfellow2014explaining}. Data augmentation manufactures variety from the examples already collected. Domain adaptation moves a model trained on one distribution to a related one, and its robotics form, domain randomization, trains in simulation with the simulation's parameters varied widely enough that reality is one more variation the system has already seen \autocite{tobin2017domain}. Ensembles trade compute for the reduction of any single model's blind spot. Regularization and human-in-the-loop review do double duty from elsewhere in this chapter.

Two limits matter more than the list. Each technique closes one route by which a system fails outside its training conditions, and none addresses the general fact of being outside them — a distinction that gets lost when robustness is reported as a single number. And all of it assumes an environment that shifts without intending to. Section~\ref{sec:6.4.4} covers the case where the shift is arranged by a state that has studied the system, which none of these techniques was built for.
